\section{Ablation Studies}
\label{sec:ablation}

\subsection{Why Hybrid Fusion is Necessary}

Na\"ive iterative refinement (feeding full conditioned output back without decoupling) degrades translation while improving rotation:

\begin{table}[h]
\centering
\caption{Na\"ive vs.\ hybrid iterative refinement over 5 iterations.}
\label{tab:naive_vs_hybrid}
\begin{tabular}{@{}lcccc@{}}
\toprule
Method & ATE$\downarrow$ & ARE$\downarrow$ & $\Delta$ATE & $\Delta$ARE \\
\midrule
Unconditioned baseline & 0.161 & 4.39° & --- & --- \\
Na\"ive self-refine (iter 5) & 0.195 & 3.60° & +21\% & $-$18\% \\
Hybrid MC (iter 2) & 0.160 & 3.13° & $-$1\% & $-$29\% \\
\bottomrule
\end{tabular}
\end{table}

The Sim(3) scale ambiguity in Pi3X outputs means that conditioned translations accumulate bias across iterations. The hybrid approach avoids this by anchoring translations to the unconditioned source.

\subsection{Noise Robustness of Pose Conditioning}

We test Pi3X's conditioning with ground truth poses corrupted by varying levels of Gaussian noise:

\begin{table}[h]
\centering
\caption{Effect of conditioning pose quality on Pi3X output (after full pipeline). Translation noise is isotropic Gaussian in the scene coordinate system.}
\label{tab:noise_robustness}
\begin{tabular}{@{}cccc@{}}
\toprule
Rot.\ noise (°) & Trans.\ noise & ATE$\downarrow$ & ARE$\downarrow$ \\
\midrule
0 (perfect GT) & 0 & 0.097 & 1.22° \\
1 & 0.05 & 0.087 & 1.30° \\
2 & 0.10 & 0.080 & 1.60° \\
5 & 0.20 & 0.077 & 3.34° \\
10 & 0.50 & 0.368 & 180° \\
20 & 1.00 & 0.951 & 180° \\
\bottomrule
\end{tabular}
\end{table}

Key observations:
\begin{itemize}
    \item Conditioning tolerates up to $\sim$5° of rotation noise
    \item Beyond 10°, the conditioning signal becomes misleading and breaks the model
    \item Our self-conditioned prior ($\sim$4.4° error) is at the edge of the useful regime, explaining the diminishing returns beyond 2--3 iterations
\end{itemize}

\subsection{Convergence Behavior}

The iterated hybrid MC method converges within 2--3 iterations:

\begin{table}[h]
\centering
\caption{Convergence of iterated hybrid MC refinement (4 samples per iteration).}
\label{tab:convergence}
\begin{tabular}{@{}ccccc@{}}
\toprule
Iteration & ATE & ARE (°) & RPE rot (°) & RPE trans \\
\midrule
0 (unconditioned) & 0.161 & 4.39 & 6.26 & 0.375 \\
1 & 0.162 & 3.43 & 4.88 & 0.321 \\
2 & 0.160 & 3.13 & 4.58 & 0.309 \\
3 & 0.163 & 3.38 & 4.70 & 0.307 \\
4 & 0.160 & 180° & 4.78 & 8.55 \\
\bottomrule
\end{tabular}
\end{table}

Iteration 2 is the sweet spot. Beyond iteration 3, the conditioned rotations begin to drift and can cause orientation flips (iter 4). We recommend stopping after 2--3 iterations or when ARE improvement is $<0.1°$.

\subsection{Sim(3) Constraint Does Not Help}

We tested Sim(3)-aligning the conditioned output back to the conditioning frame before each iteration (to prevent scale drift). Results are identical to na\"ive iteration, confirming that translation degradation is not caused by scale drift but by systematic bias in how Pi3X uses pose conditioning:

\begin{table}[h]
\centering
\caption{Sim(3)-constrained vs.\ na\"ive iteration (5 iterations).}
\label{tab:sim3}
\begin{tabular}{@{}ccccc@{}}
\toprule
Iter & Sim3 ATE & Sim3 ARE & Na\"ive ATE & Na\"ive ARE \\
\midrule
0 & 0.152 & 4.69° & 0.152 & 4.69° \\
1 & 0.160 & 3.50° & 0.160 & 3.50° \\
2 & 0.180 & 3.31° & 0.179 & 3.29° \\
5 & 0.197 & 3.61° & 0.195 & 3.59° \\
\bottomrule
\end{tabular}
\end{table}
